\ticket{2}{Принципы ООП и обобщенное программирование}

\begin{examplan}
    \item Дать определение объектно-ориентированного программирования.
    \item Раскрыть абстракцию, инкапсуляцию и полиморфизм.
    \item Перечислить виды полиморфизма и отличить статический полиморфизм от динамического.
    \item Объяснить концепцию обобщенного программирования и его связь с параметрическим полиморфизмом.
\end{examplan}

\subsection*{Короткая версия для письменного ответа}

\textbf{Объектно-ориентированное программирование} --- это парадигма, в которой программа представляется как совокупность взаимодействующих объектов. Объект имеет состояние, поведение и идентичность; обычно он является экземпляром класса. Класс описывает общие атрибуты и методы для группы объектов. Классы могут образовывать иерархии наследования, но ООП не сводится только к наследованию.

Главная идея ООП --- моделировать предметную область через объекты и их взаимодействия. Это помогает управлять сложностью, локализовать изменения, повторно использовать код и строить расширяемые системы.

\subsection*{Абстракция}

\textbf{Абстракция} --- это выделение существенных свойств и поведения объекта при отвлечении от несущественных деталей реализации. Абстракция отвечает на вопрос: \textit{что важно знать об объекте в данной задаче}.

\textbf{Зачем нужна абстракция:}
\begin{itemize}
    \item уменьшает сложность системы;
    \item позволяет работать с объектом через понятный интерфейс;
    \item отделяет смысловую модель предметной области от технических деталей реализации.
\end{itemize}

\textbf{Как реализуется:}
\begin{itemize}
    \item через классы, которые задают общие поля и методы объектов;
    \item через интерфейсы и абстрактные классы, которые описывают контракт: что объект должен уметь делать, не фиксируя конкретный способ реализации.
\end{itemize}

\textbf{Пример.} В библиотечной системе сущность \texttt{Book} можно описать через название, автора, ISBN и операции \texttt{issue()} и \texttt{returnBook()}. Детали хранения в базе данных или расположения книги на полке не входят в эту абстракцию, если они не важны для текущей задачи.

\subsection*{Инкапсуляция}

\textbf{Инкапсуляция} --- это объединение данных и методов, работающих с этими данными, внутри объекта, а также ограничение прямого доступа к внутреннему состоянию объекта.

Инкапсуляция тесно связана с \textbf{сокрытием данных}: объект предоставляет внешний интерфейс, а детали реализации оставляет внутри. Поэтому другие части программы зависят не от внутреннего устройства объекта, а от его публичного поведения.

\textbf{Зачем нужна инкапсуляция:}
\begin{itemize}
    \item защищает состояние объекта от некорректного изменения;
    \item сохраняет инварианты класса;
    \item повышает модульность: реализацию можно менять, если публичный интерфейс остается прежним;
    \item упрощает сопровождение и тестирование.
\end{itemize}

\textbf{Как реализуется:}
\begin{itemize}
    \item модификаторами доступа: \texttt{private}, \texttt{protected}, \texttt{public};
    \item публичными методами, через которые объект контролирует чтение и изменение своего состояния;
    \item методами с проверками, а не прямым доступом к полям.
\end{itemize}

\textbf{Пример.} В классе банковского счета поле \texttt{balance} делают закрытым. Изменять баланс можно только через методы \texttt{deposit(amount)} и \texttt{withdraw(amount)}, где проверяется положительность суммы и достаточность средств.

\subsection*{Полиморфизм}

\textbf{Полиморфизм} --- это возможность использовать единый интерфейс для работы с объектами или значениями разных типов, при этом конкретное поведение определяется типом объекта или способом параметризации.

В ООП полиморфизм позволяет писать код, который зависит от общего интерфейса, а не от конкретных классов. Это уменьшает число проверок вида \texttt{if}/\texttt{switch} по типу объекта и облегчает расширение системы.

\subsection*{Виды полиморфизма}

\textbf{1. Полиморфизм подтипов} также называют включающим полиморфизмом. Объект класса-потомка можно использовать там, где ожидается объект базового класса или интерфейса. Конкретная реализация метода выбирается во время выполнения через динамическую диспетчеризацию.

\begin{verbatim}
Shape[] shapes = { new Circle(), new Rectangle() };
for (Shape s : shapes) {
    s.draw(); // вызывается draw конкретной фигуры
}
\end{verbatim}

Этот вид полиморфизма связан с наследованием, интерфейсами, переопределением методов и поздним связыванием.

\textbf{2. Ad-hoc полиморфизм} означает, что одно имя операции используется для разных типов, но реализации могут быть независимыми. Типичные формы:
\begin{itemize}
    \item перегрузка функций или методов;
    \item перегрузка операторов.
\end{itemize}

Например, метод \texttt{print} может иметь версии для \texttt{int}, \texttt{double} и \texttt{String}. Обычно выбор версии происходит на этапе компиляции, поэтому это статический полиморфизм.

\textbf{3. Параметрический полиморфизм} позволяет писать функции и классы с параметрами типов. Один алгоритм работает для разных типов, если эти типы удовлетворяют нужным требованиям.

Примеры: \texttt{List<T>} в Java/C\#, \texttt{vector<T>} и шаблонные функции в C++.

\textbf{4. Полиморфизм приведения} связан с автоматическими или явными преобразованиями типов. Например, значение типа \texttt{int} может быть приведено к \texttt{double}. Этот вид обычно менее важен для ООП, но его иногда выделяют в общей классификации полиморфизма.

\subsection*{Статический и динамический полиморфизм}

\begin{description}
    \item[Статический полиморфизм] разрешается на этапе компиляции. К нему относят перегрузку и многие реализации обобщений/шаблонов.
    \item[Динамический полиморфизм] разрешается во время выполнения. К нему относят подтипный полиморфизм с виртуальными методами.
\end{description}

\subsection*{Обобщенное программирование}

\textbf{Обобщенное программирование} --- это подход, при котором алгоритмы и структуры данных пишутся независимо от конкретных типов данных. Типы задаются как параметры, а один и тот же код можно использовать для разных типов.

\textbf{Цели обобщенного программирования:}
\begin{itemize}
    \item повторное использование кода;
    \item типобезопасность: ошибки типов обнаруживаются как можно раньше;
    \item отделение алгоритма от конкретного представления данных;
    \item эффективность, особенно в языках, где обобщенный код специализируется компилятором.
\end{itemize}

\textbf{Реализация в языках:}
\begin{itemize}
    \item \textbf{C++}: шаблоны \texttt{template}, стандартная библиотека STL, контейнеры \texttt{vector<T>}, \texttt{list<T>}, алгоритмы вроде \texttt{sort};
    \item \textbf{Java}: generics, например \texttt{ArrayList<E>} и \texttt{HashMap<K,V>}; в Java используется стирание типов;
    \item \textbf{C\#}: generics, например \texttt{List<T>} и \texttt{Dictionary<TKey,TValue>}; информация о типах сохраняется во время выполнения.
\end{itemize}

Обобщенное программирование практически реализует идею \textbf{параметрического полиморфизма}: одна функция или структура данных параметризуется типом. В отличие от подтипного полиморфизма, здесь не обязательно нужна иерархия наследования. В отличие от перегрузки, обычно пишется одна общая реализация, а не набор отдельных реализаций для каждого типа.

\subsection*{Сравнение ключевых понятий}

\begin{description}
    \item[Абстракция] скрывает несущественные детали и выделяет важную модель объекта.
    \item[Инкапсуляция] скрывает внутреннее состояние и защищает его через контролируемый интерфейс.
    \item[Полиморфизм] позволяет работать с разными типами через общий интерфейс или общий код.
    \item[Обобщенное программирование] позволяет писать алгоритмы и структуры данных, параметризованные типами.
\end{description}

\begin{important}
    \item ООП организует программу вокруг объектов, их состояния, поведения и взаимодействий.
    \item Абстракция отвечает за выделение существенного, инкапсуляция --- за защиту внутреннего состояния, полиморфизм --- за единый интерфейс для разных реализаций.
    \item Подтипный полиморфизм обычно динамический: выбор метода происходит во время выполнения.
    \item Перегрузка и шаблоны/дженерики чаще относятся к статическому полиморфизму.
    \item Обобщенное программирование основано на параметрическом полиморфизме и позволяет не переписывать алгоритм под каждый тип.
\end{important}

